\documentclass[12pt]{article}
\usepackage[utf8]{inputenc}
\usepackage[T1]{fontenc}
\usepackage{hyperref}
\title{Tocho Language Resources for K-12 Curriculum Development}
\author{Compiled by Claude Sonnet 5 (claude-sonnet-5) for Andr\'as Kornai}
\date{}
\begin{document}
\maketitle

\section{Introduction}
Tocho (ISO 639-3: taz; Glottocode: toch1257) is a Niger--Congo language of the Talodi group, spoken in South Kordofan, Sudan. Wikipedia rates it Severely Endangered and gives about 2,700 speakers (2013); other sources give about 2,650 or about 3,800. Very little is published about Tocho. No school materials, primers or native-language texts turned up in web search, and no evidence was found that it is taught in schools.

\section{Primary Educational Resources}
No primers, textbooks or school programmes for Tocho were found in search results.

\section{Linguistic and Reference Materials}
\begin{itemize}
\item Wikipedia, ``Tocho language'': overview of the phonology and noun-class prefixes. It draws on Alaki and Norton (2013), with data collected in Khartoum, so it should be checked with village speakers. \url{https://en.wikipedia.org/wiki/Tocho_language}
\item Max Planck EVA, Tocho numerals (data from Russell Norton): a base-10 system with special forms for 6--9. \url{https://lingweb.eva.mpg.de/channumerals/Tocho.htm}
\item Omniglot, Tocho: Latin script, speaker locations in South Kordofan. \url{https://www.omniglot.com/writing/tocho.htm}
\item Endangered Languages Project entry for Tocho (Talodi, Moro Hills area). \url{https://endangeredlanguages.com/lang/5115}
\end{itemize}
A search snippet also pointed to Schadeberg, \emph{A Survey of Kordofanian, Vol.\ 2: The Talodi Group} (1981) as the main older reference. I did not find an online copy or URL for it.

\section{Religious and Community Resources}
\begin{itemize}
\item Joshua Project, Tocho in Sudan: people-group profile. It says scripture translation has started but no Bible is available. Its page carries an ``AI generated'' note and its population figures look like whole-people counts, so treat it cautiously. \url{https://joshuaproject.net/people_groups/13818}
\end{itemize}

\section{Digital and Audio Resources}
No audio, video, corpus or software resources for Tocho were found.

\section{Curriculum Development Context}
Tocho has no known written educational tradition, and most children's schooling in Sudan is in Arabic. A native-language curriculum would first need community-led orthography work, documentation and teacher materials, starting from the linguistic sources above and fieldwork with speakers. Security and displacement in South Kordofan may complicate fieldwork.
\end{document}